\documentclass{clean_cv}
% Add a BibTeX-style file encoding all of your publications to include here. You can export this from Zotero. Only include
% publications you want to appear here!
\addbibresource{publications.bib}
\author{\textbf{Matthew Prest}}
\begin{document}

\maketitle
% In this section, you can use any of the FontAwesome icons. The commands \faCenter and \faCenterStyle have been defined to properly center the icons
% when using the default font settings.
%
% You can use any of the icons listed in the fontawesome5 package documentation (https://ctan.math.utah.edu/ctan/tex-archive/fonts/fontawesome5/doc/fontawesome5.pdf)
% If you need to specify a specific style (as is done here for the address card), you should use the two-argument \faCenterCycle command
\begin{center}
\begin{tabular}{lll}
    \faCenter{globe} \href{https://www.mtprest.com/}{{mtprest.com}}
    \faCenter{envelope} \href{mailto:MatthewPrestNZ@gmail.com}{{MatthewPrestNZ@gmail.com}}  & \faCenter{linkedin} \href{https://www.linkedin.com/in/matthew-prest-4165aa1a7/}{{Matthew Prest}} & \faCenter{github} \href{https://github.com/mtp354}{{mtp354}} \faCenter{phone-alt} 917-302-6890
\end{tabular}

\end{center}


\section{Education \& Certification}

\begin{datetabular}{7em}

\dateentry{2025}{
\textbf{IBM Quantum Computation using Qiskit -- }
\textit{Certified Associate Developer}
\eatvspace}

\dateentry{2024--Present}{
\textbf{City University of New York -- }
\textit{Ph.D. in Physics}
\begin{itemize}
    \item Research assistant working in quantum cryptography and information theory
    \item Coursework focused on quantum algorithms, quantum networks, entanglement distribution, and hybrid classical-quantum systems.
    \item External reviewer under supervision for \textit{Information and Computation} and \textit{Crypto2026}
\end{itemize}\eatvspace}

\dateentry{2022}{
\textbf{DataCamp Data Scientist Professional License -- }
\textit{Certified Python Developer}
\eatvspace}

\dateentry{2019--2021}{
\textbf{New York University -- }
\textit{Master of Science (M.S.) in Physics}
\begin{itemize}
    \item Thesis: Quantum Computing using Cold Atomic Ensembles
    \item NYU Shanghai Fellowship Recipient
    \item Coursework in neutral atom systems and Noisy Intermediate-Scale Quantum Computing
\end{itemize}
\eatvspace}
\dateentry{2015--2019}{
\textbf{St. Louis University -- }
\textit{Bachelor of Science (B.S.) in Physics \& Mathematics | cum laude}
\begin{itemize}
    \item Graduation Speaker
    \item All-Conference Division I Student-Athlete
    \item Pi Mu Epsilon National Honorary Mathematics Society
\end{itemize}
\eatvspace}


\end{datetabular}


\section{Relevant Experience}

\begin{datetabular}{7em}
\dateentry{2025}{
\textbf{Qiskit Community Advocate -- IBM}
\begin{itemize}
\item Promoted Qiskit tools to campus researchers and participated in IBM Quantum/Qiskit seminars.
\item Currently developing open-source Qiskit contributions for efficient modular-reduction circuits.
\end{itemize}\eatvspace}

\end{datetabular}

\begin{datetabular}{7em}
\dateentry{2024--2025}{
\textbf{Research Assistant -- CUNY}
\begin{itemize}
\item Designed fault-tolerant quantum cryptographic protocols for distributed quantum computing.
\item Developed and implemented a protocol for restricted state quantum broadcasting.
\item Applied iterative numerical methods in Python to simulate noisy and disordered many-body interactions in low-dimensional materials.

\end{itemize}\eatvspace}
% The command \eatvspace should be used if extra space appears at the end of a datetabular environment.
\end{datetabular}

\begin{datetabular}{7em}
\dateentry{2021--2024}{
\textbf{Programmer/Statistician -- Columbia University}
\begin{itemize}
\item Exceeded expectations in every performance review over 3 years under multiple managers.
\item Implemented stochastic optimization algorithms in Python, including simulated annealing and gradient descent, using cloud-based high-performance computing resources to improve cancer-forecasting accuracy. This work supported multiple conference presentations to audiences of more than 100.
\item Created a version-controlled Python/R library, hosted on GitHub, to simulate and forecast cancer trends and analyze screening/intervention impacts; the library is now used by multiple research groups.
\end{itemize}\eatvspace}
\end{datetabular}

\begin{datetabular}{7em}
\dateentry{2019--2021}{
\textbf{Teaching/Research Assistant -- New York University}
\begin{itemize}
\item Simulated and analyzed a quantum-entanglement scheme, resulting in a peer-reviewed publication.
\item Designed and constructed laboratory equipment used in multiple subsequent projects.
\item Taught multiple class sections with over 20 students each; designed quizzes, graded assignments, and analyzed educational-outcomes data.
\end{itemize}\eatvspace}
% The command \eatvspace should be used if extra space appears at the end of a datetabular environment.
\end{datetabular}
\eatvspace
\section{Skills}
\textbf{Scientific Programming:}
Python, PyTorch, Qiskit, C++, R, SQL, Mathematica, MATLAB, LaTeX\\
\textbf{Software Practices:}
Git/GitHub, open-source development, version control, unit testing\\
\textbf{Spoken Languages:}
English (Native), Mandarin (Limited) \\



\section{Selected Publications \& Conference Proceedings}
\nocite{*} % Loads every entry from the attached .bib file
% Highlight takes three entries, given name, given name initials, and family name.
% If you have a middle initial, this call looks like:
% \highlightauthorname{Bob H.}{B. H.}{Smith}
\highlightauthorname{Matthew}{M.}{Prest} 
\begin{datetabular}{7em}
%printbibyear has been defined to only print entries from a given citation year.

\dateentry{Quantum}{\printbibyear{0001}}
\dateentry{Other}{\printbibyear{0002}}


\end{datetabular}


\section{Other Experience}
\begin{datetabular}{7em}
\dateentry{2025}{
\textbf{QuEra Quantum Creators Hackathon}
\begin{itemize}
\item Led a student team in building an educational animation explaining the hardware components of neutral-atom quantum computers.
\item Worked with QuEra staff to communicate technical quantum-computing concepts to a non-specialist creative audience.
\end{itemize}\eatvspace}
% The command \eatvspace should be used if extra space appears at the end of a datetabular environment.
\end{datetabular}

\begin{datetabular}{7em}
\dateentry{2019}{
\textbf{Intern -- Bluebox Technical Ltd.}
\begin{itemize}
\item Performed import quality control of industrial polymer-based adhesives using reactant analysis.
\item Analyzed local market demand for product sizing to optimize imports.
\item Presented technical information to non-technical stakeholders to support decision-making.
\end{itemize}\eatvspace}
% The command \eatvspace should be used if extra space appears at the end of a datetabular environment.
\end{datetabular}

\begin{datetabular}{7em}
\dateentry{2017--2019}{
\textbf{Undergraduate Research Assistant -- Saint Louis University}
\begin{itemize}
\item Used LabVIEW to integrate oscilloscope data into existing data-acquisition systems.
\item Operated a chemical vapor deposition (CVD) system to grow carbon nanotubes and graphene for tunable microwave resonator devices.
\item Used C++ to control a vector network analyzer and measure microwave-resonator response as a function of bias current.
\item Designed and constructed bias-current lines for an adiabatic demagnetization refrigerator.
\end{itemize}\eatvspace}
% The command \eatvspace should be used if extra space appears at the end of a datetabular environment.
\end{datetabular}

\begin{datetabular}{7em}
\dateentry{2017--2024}{
\textbf{Academic Tutor}
\begin{itemize}
\item Tutored high-school students in algebra, geometry, calculus, and SAT preparation.
\item Tutored undergraduate students in calculus, differential equations, linear algebra, object-oriented programming, real analysis, machine learning, thermodynamics, and statistical mechanics; also provided GRE and MCAT preparation.

\end{itemize}\eatvspace}
% The command \eatvspace should be used if extra space appears at the end of a datetabular environment.
\end{datetabular}






\end{document}
